% 本文件由 LaTeX 排版 Agent 根据有效 translation.json 自动生成。
% author=Athenagoras; work_id=0206; title=论死者的复活

\chapter{论死者的复活}

% structure: p0001 source=h1 target=omitted reason=page-title-already-rendered-by-parent

% structure: p0002 source=h2 target=section reason=relative-heading-rank; base=h2

\section{第一章 为真理辩护应先于对真理的讨论}

对于每一种与事物之真理相符的意见和道理，都会滋生某种虚假；而虚假之所以出现，并非因为它从某一根本原则或某项针对当前问题的特殊原因中自然产生，而是因为它是由那些重视虚假种子的人故意发明的，因为这种种子有败坏真理的倾向。这一点，首先从昔时那些专注于此等探究的人，以及他们与其前人、同辈的不一致，便显而易见；其次，从时下所盛行讨论的混乱状态，也同样昭然若揭。因为这样的人，没有让任何真理免受他们的诽谤攻击——无论是天主的存在，祂的知识，祂的作为，还是那些循规蹈矩严格由此而来，为我们将虔敬道理描绘出来的书籍，概莫能外。相反，他们中有些人完全且一劳永逸地绝望地放弃了有关这些事的真理，有些人则歪曲真理以适合己见，还有些人甚至对显而易见的事故意怀疑。因此，我认为那些关注此类问题的人应采取两种论证路线，一为真理辩护，一为探讨真理：为真理辩护，是为不信者和怀疑者；探讨真理，是为那些坦诚且愿意接受真理的人。是故，凡愿研究这些事的人，务要留意每一事例的必要性所要求的，并据此调整他们的讨论；要使他们处理这些主题的次序适用于恰当的时机，而不可为了显得总是保持同一方法，却忽视合适性及各主题所应占的位置。因为，就证明和自然次序而言，关于真理的论文总优先于为真理辩护的论文；但为了更大的益处，次序必须颠倒过来，为真理辩护的论证应先于关于真理的论证。农夫若不先清除野草树木，以及一切有害于良种的杂物，便不能恰当地将种子撒入土中；医生若未先除去体内的疾病，或阻止其逼近，便也不能将有益的药物引入需要照护的身体。同样，欲教导真理的人，如果听者心中仍潜伏着谬见，阻塞其论证的进入，那么无论他如何讲说，也必不能说服任何人。因此，为了更大的益处，我本人有时也将为真理辩护的论证置于关于真理的论证之前；而在目前，就此事的要求而言，我认为在处理复活问题时遵循同一方法并非无益。因为关于这个主题，我们同样发现有些人完全不信，有些人则心存怀疑，甚至在已接受了首要原则的人中，也有些人与怀疑者一样茫然不知何所信从；而最不可理喻的是，他们之所以处于这种心态，并非由于这些事情本身有丝毫根据使他们不信，也未能为他们的不信或困惑找到任何合理的理由。\par

% structure: p0004 source=h2 target=section reason=relative-heading-rank; base=h2

\section{第二章 复活并非不可能}

让我们按照我指出的方式来考虑这个主题。如果所有的不信并非出于轻率和欠考虑，而是基于强有力的理由，并伴随着属于真理的确定性而产生[那很好]；因为当他们的不信所关乎的事物本身，在他们看来不值得相信时，不信便维持着公正的表象；但不信那些不值得不信的事物，乃是那些对真理没有作出健全判断之人的行为。因此，那些不信或怀疑\OriginalTerm{复活}{resurrection}的人，理当形成他们的观点，不是根据任何草率形成的看法，也不是根据放荡之人所接受的，而是或者不将人的起源归于任何原因（这种观念极易被驳倒），或者，将万物的原因归于天主，并牢牢记住这一信条所包含的原则，由此证明\OriginalTerm{复活}{resurrection}完全不值得相信。如果他们能够证明，天主要么不可能，要么违背祂的旨意，将死了的，或甚至完全分解为\OriginalTerm{元素}{elements}的身体，重新组合并聚集起来，以构成同样的人，那他们便能成功。如果他们不能，就当停止这种不敬虔的不信，以及对神圣事物的亵渎：因为，他们声称那是不可能的，或不合乎天主圣意的说法，并非真理，这在我接下来所要说的，将清楚地显明。严格来说，一件事对一个人被认为是不可能，当这件事属于这样一种情况：他要么不知道当作什么，要么没有足够的能力去适当地完成那已知的事。因为，凡对任何需要做的事无知的人，完全不能尝试或去做他所无知的事；同样，极为清楚知道要做什么，以及藉着什么方法，如何去做，但要么完全没有能力去做那已知的事，要么能力不足，如果他明智且考虑自己的能力，他甚至不会尝试；而如果他未经充分考虑就去尝试，他也无法达成目的。但是，天主不可能无知：无论是关于将要复活的身体的性质，涉及到完整的肢体和构成它们的\OriginalTerm{微粒}{particles}，还是关于每个分解的微粒去了何处，\OriginalTerm{元素}{elements}的哪一部分接纳了那分解的，并与它亲和的物质结合了——尽管在人类看来，那已按本性重新与\OriginalTerm{宇宙}{universe}结合的，再次从中分离出来，似乎是完全不可能。因为，在每个人的独特形成之前，祂既未隐瞒构成人类身体的\OriginalTerm{元素}{elements}的特性，也未隐瞒祂将从其中取用那些似乎适合于形成人类身体的那些部分的特性；那么，显而易见，在整体\OriginalTerm{分解}{dissolution}之后，祂也不会无知，祂为构成每个身体所取用的每个微粒去了何方。因为，相对于我们中间现今存在的秩序，以及我们对其他事物所形成的判断，预知尚未发生的事是更伟大的；但是，相对于天主的威严和智慧，二者皆是合乎本性的，预知尚未存在的事，与知晓已分解的事，同样是容易的。\par

% structure: p0006 source=h2 target=section reason=relative-heading-rank; base=h2

\section{第3章 能创造的，也能使死人复活}

此外，祂的能力足以使死去的身体\OriginalTerm{复活}{raising}，这由这些身体本身的创造所表明。因为，如果当它们不存在时，祂在最初形成时造了人的身体，以及它们的\OriginalTerm{原始元素}{original elements}，那么，当它们\OriginalTerm{分解}{dissolved}时，无论以何种方式发生，祂都会以同样的容易再次将它们复活：因为这对祂同样是可能的。并且，这并不损害论证，即便有些人假设\OriginalTerm{最初的起始}{first beginnings}来自\OriginalTerm{物质}{matter}，或者至少人的身体来源于那些作为最初材料的\OriginalTerm{元素}{elements}，或是来源于\OriginalTerm{种子}{seed}。因为那种能力能够给它们视为\OriginalTerm{无形物质}{shapeless matter}的东西赋予形状，当它缺乏形式和秩序时，用许多不同的形式装饰它，将\OriginalTerm{元素}{elements}的各个部分聚集为一，将本是单一简单的\OriginalTerm{种子}{seed}分成许多，\OriginalTerm{组织}{organize}那未经组织的，并将生命赋予那没有生命的——同一种能力能够将分解的重新结合，将\OriginalTerm{倒下的}{prostrate}扶起，使死人恢复生命，使\OriginalTerm{可朽坏的}{corruptible}进入\OriginalTerm{不朽坏}{incorruption}的状态。并且，同一位存在者，同一种能力和技巧，将能够把那已经破碎并分布到多种动物之中的部分分离出来——那些动物习惯于求助于这样的身体，并以它们来满足食欲——我要说，将这部分分离出来，并再次与适当的肢体和肢体的部分结合，无论它是进入了那些动物中的某一个，还是进入了多个，或从那里又进入其他动物，或者，在与这些动物一同分解之后，根据自然法则，又回到了\OriginalTerm{原始元素}{original elements}，被分解为这些元素——这件事似乎极大地困扰了一些人，甚至包括那些因智慧而受钦佩的人，他们，我不知道为何，认为这些由群众提出的怀疑值得认真关注。\par

% structure: p0008 source=h2 target=section reason=relative-heading-rank; base=h2

\section{第4章 来自某些人体成为他人一部分之事实的异议}

这些人说，许多死于船只失事与河流中者的身体成了鱼的食物，许多战死者，或因其他不幸原因或境况未被埋葬者，暴露在外，成为碰巧遇见的任何动物的食物。既然身体如此被吞吃，组成身体的肢体与部分被分解并散布于众多动物中间，又藉着营养作用融入那些以它们为食的动物体内——他们首先说，将这些部分从这些动物体内分离出来是不可能的；此外，他们又提出另一更为困难的情况。当那些以人身为食、适合作为人类食物的动物经过他们的胃，并融入那些食用它们的人的身体时，他们说，人体中曾作为食物动物营养的部分，必然要进入其他人的身体，因为这些动物在此期间因它们而获得滋养，便将从所食之人得来的营养传递给那些以它们为食的人。然后，他们又悲惨地加上：在饥荒与疯狂中人们吞食自己的孩子，以及仇敌诡计下父母吃掉自己的儿女，还有那著名的\OriginalTerm{米底亚盛宴}{Median feast}，和\Person{提厄斯忒斯}的悲剧筵席；他们还加上\Place{希腊人}与蛮族中发生过的其他类似闻所未闻的事件。从这些事中，他们自以为便确立了复活的不可能性，其理由是：同样的部分不可能与这一套身体一同复活，又与另一套身体一同复活；因为要么原来占有者的身体无法重组成形，因其构成部分已转入他人体内，要么若将这些部分归还给原来者，则最后占有者的身体就会有所短缺。\par

% structure: p0010 source=h2 target=section reason=relative-heading-rank; base=h2

\section{第五章 论消化与营养的过程}

在我看来，这些人首先不知道那形成并治理此宇宙的天主的能力与技巧，祂依照每种动物的本性与种类，赋予它们合宜且相应的食物，既未命定自然界的一切都要与每一种身体联合与结合，也并非不能将如此联合之物加以分离，而是准许各个受造物的本性，去行或接受那自然适合于它的事，有时也依照祂的意愿与目的，加以阻拦、准许或禁止；此外，他们也没有思考过那些提供滋养或接受滋养的受造物各自的能力与本性。否则他们便会知道，并非所有在外部必要压力下被当作食物的东西，都成了那动物合宜的营养；反之，有些东西刚一接触到\OriginalTerm{胃的皱襞}{plicatures of the stomach}，就立即腐坏，被呕出、排泄，或以其他方式被处理掉，以致连初次而自然的消化都不曾经历片刻，更不要说融入那有待滋养之物了；同样地，也并非所有已经在胃中消化并经历了初次变化的食物，都实际到达待滋养的部分，因为有些甚至在胃中就丧失了营养能力，有些则在第二次变化期间，而在肝中进行的消化则被分离出来，转变成其他缺乏营养能力的东西；不，肝中所发生的变化，并不全都成为人的营养，所变之物依照其本然的用意，作为废物被分离；而留在需要滋养的肢体与部分本身中的营养，有时也转变成其他东西，这取决于那多或少、且易于腐化或将近旁之物转化为自己的因素是否占优势。\par

% structure: p0012 source=h2 target=section reason=relative-heading-rank; base=h2

\section{第六章 一切无用或有害之物皆被弃绝}

既然，在所有动物中，其本性存在着巨大的差异，而与本性相符的营养本身也因每种动物及被滋养的身体而异；并且在每种动物的营养过程中，存在三重洁净与分离，由此可知，凡与该动物营养相异的，必然被完全摧毁并带至其自然之处，或转化为他物，因为它不能与之融合；滋养之物的能力必须适合于被滋养动物的本性，并与该动物的能力相符；而这种能力，在通过为此目的而设置的过滤筛滤，并经过自然净化方式的彻底净化之后，必须成为对质料最真确的增添——事实上，这乃是任何以正确名称称呼事物的人会称为真正营养的唯一东西；因为它排斥一切对被滋养动物的体质有碍或有害的外物，以及仅仅为了填饱胃囊、满足口腹之欲而引入的大量多余食物。无人能怀疑，这种营养会与受滋养的身体结合，交织并融合于所有的肢体和肢体的部分之中；但那些与本性格格不入且相反的东西，若与一种更强的能力接触，便会迅速败坏，而它却容易摧毁被它克服之物，并转化为有害的体液和有毒的性质，因为它不能产生任何与将要受滋养的身体同质或相合的东西。一个非常清楚的证明是，在许多被滋养的动物中，痛苦、疾病或死亡皆由此而来：若因过于贪婪的食欲，它们连同食物一起摄入了某些有毒且违反本性的东西；这当然会导致被滋养的身体彻底毁灭，因为受滋养之物是由与其同质且符合其本性的物质所滋养，却被相反的东西摧毁。因此，如果按照动物本性的不同，已预备了适合其本性的各种食物，而动物所摄取的东西中，没有一样——甚至连其偶然的部分也不容许——能与受滋养的身体融合，只有那部分经过了完全消化而净化，并为与特定身体结合而经历了彻底变化，并适应了将接受营养的部位的东西——那么，非常清楚，任何违反本性的东西都无法与那些身体结合，因为它们不合适也不相匹配的营养，要么在产生某种粗劣败坏的体液之前就通过肠道排出；要么，如果停留时间更长，就会引发难以治愈的痛苦或疾病，同时摧毁原有的自然营养，甚至需要营养的肉体本身。但即使它最终被某些药物、更好的食物或自然力量所克服而排出，不造成显著伤害也是无法摆脱的，因为它对自然之物不显丝毫平和之态，这乃是因为它不能与本性和谐共存。\par

% structure: p0014 source=h2 target=section reason=relative-heading-rank; base=h2

\section{第七章 复活的身体与现今的不同}

不，即使我们假定来自这些东西的营养（姑且如此称呼，以更符合通常的说法），虽然违反本性，却仍被分离并转化为身体所含的某种湿、或干、或温、或冷的物质，我们的对手也绝不会因这一让步而获益：因为复活的身体是由那些原本属于它们自己的部分重新构成的，而上述事物中没有一样是这样的部分，也没有部位的形式或位置；不，它并不始终与受滋养的身体部分同在，也不与复活的部分一同复活，因为血液、黏液、胆汁或气息不再对生命有任何贡献。再者，那时受滋养的身体也不再需要它们曾经需要的东西，因为随着受滋养身体的匮乏与败坏，对于那些曾用来滋养它们的东西的需求也已消失。此外还必须补充说，如果我们假设由这类营养所产生的变化可远达于肌肉，那么在这种情况下，也没有必要认为，最近被那种食物改变过的肌肉，如果结合到另一个人的身体上，就应再次作为一个部分参与那身体的构成，因为吸收它的肌肉并不总是保留它所吸收的东西，这样融入的肌肉也不会恒久与它所附加之处共存，而是会经历多种多样的变化：时而因辛劳或忧虑而消散，时而因悲伤、烦扰或疾病而耗损，并因受热或受寒引起的失调而受影响，这些与肌肉和脂肪一同变化的体液并未接受营养以保持其原状。但是，肌肉既然会遭此变化，我们就应发现，被不适合的食物所滋养的肌肉，其变化程度要大得多；时而因所接受之物而肿胀、发胖，然后又以某种方式将其排出，体积减小，这可以由前面提到的一个或多个原因造成；而留存在各部分的，唯有那适合于连结、覆盖或温暖由本性所选定的肌肉，并附着于那些部分，借以维系依循本性的生命，并履行那生命之辛劳的东西。因此，无论我们刚才所进行的探究得到公正的评判，还是对我们立场的反对意见被认可，在任何一种情况下，都不能证明对手所言为真，而且人的身体也绝不可能与那些本性相同的身体结合，无论是在什么时候，人因无知、被他人误导而吃了这样的身体，还是自己出于欲望或疯狂，用形体相似者的身体玷污了自己；因为我们非常清楚，有些野兽具有人形，或者是一种由人与野兽混合而成的本性，正如较为大胆的诗人们惯于描绘的那样。\par

% structure: p0016 source=h2 target=section reason=relative-heading-rank; base=h2

\section{第八章 人的肉并非人类合宜或天然的食物}

但是，对于那些并非注定成为任何动物之食物、而只是为了顺应自然之荣耀被埋入土中的身体，又何必多谈呢？因为天主并没有将任何一种动物分配给同类作为食物，尽管某些异类动物按自然天性可作为食物。如果，确实，他们能证明人肉被分配给人作食物，那么他们彼此相食便是顺乎自然，就如同任何其他自然容许之事一样，并且那些胆敢如此说的人也就可以把至亲好友的身体当作美味佳肴，尽情享用，仿佛这特别适合他们，甚至用同样的菜肴款待他们的活人朋友。但是，如果连谈论此事都是不合法的，如果人吃人肉是最可憎恶、最令人厌恶的，比任何其他不合法、不自然的食物或行为更为可厌；如果违反自然的东西绝不可能转化为需要滋养的四肢和部位的营养，而不能成为营养的东西绝对无法与它本不适于滋养的身体相结合——那么，人的身体就绝不可能与类似的身体结合，因为这种滋养是违反自然的，即使由于某种极为惨痛的不幸，这身体多次通过他们的胃；相反，当脱离了滋养力的影响，并再次散入那些最初获得起源的宇宙各部分时，它们便与这些部分结合，时间长短各有定数；然后，凭那规定了一切动物本性并赋予各自特有能力的天主的技巧与大能，它们再被分离出来，彼此适当地结合，无论它们是被火烧毁，被水腐烂，被野兽或任何其他动物吞食，还是在其他部分之前从整个身体分离并消散；而它们重新彼此联合，占据同一位置，为的是精确构建和形成的同一身体，以及那已死、甚至完全消散者的复活与生命。然而，对这些话题进一步详述并不合适；因为所有的人——至少那些不是半野兽的人——都对它们得出了相同的结论。\par

% structure: p0018 source=h2 target=section reason=relative-heading-rank; base=h2

\section{第九章 以人的无能论证之荒谬}

鉴于有许多更为重要的事项有待我们当前的探究，我请求暂免答复那些诉诸人的作品、甚至其制造者的人——他们无法将那些破碎、因时日而磨损、或遭其他毁坏的作品重新制造出来，进而以陶匠与木匠的类比试图说明：天主既不能愿望、即便愿望了也无能力使已死或已消散的身体复活——他们未曾考虑，通过这样的推理，他们是在极度地侮辱天主，因为他们将截然不同之物的能力，或更确切地说，将这些物品使用者的本性，等同起来，并且将艺术作品与自然之作相比较。对这样的论点给予任何严肃的注意，并不免于非难，因为回答肤浅而琐屑的反驳其实是愚蠢的。可以肯定，说在人不可能的事，在天主却是可能的，这不但更为合理，而且绝对真实。而且，如果这个陈述本身是可能的，并且通过我们刚才所进行的全部探究，理性也表明这事是可能的，那么十分清楚，它并非不可能。不，这也并非天主所不能愿望之事。\par

% structure: p0020 source=h2 target=section reason=relative-heading-rank; base=h2

\section{第十章 无法证明天主不愿意复活}

凡与祂的旨意不符之事，要么是因为它不义，要么是因为它与祂不相称。再者，不义要么涉及那将要复活的人，要么涉及他之外的其他人。但显然，他之外、且被算作存在之物中的任何一个，都未受到伤害。精神性存在\OriginalTerm{精神性存在}{\GreekText{νοηταὶ} \GreekText{φύσεις}}不会因为人的复活而受到伤害，因为人的复活并不阻碍它们的存在，也不会给它们带来任何损失或暴力；同样，非理性或非生命之物的本性也不会遭受损害，因为它们将在复活之后不复存在，而对于不存在的东西，不可能行不义。但即使有人设想它们将永远存在，它们也不会因为人类身体的更新而遭受损害：因为，如果现在，它们在服务于人的本性及其需求而为人所需时，承受着轭和各种苦役，并未遭受不义，那么，当人成为不朽且无所缺乏，不再需要它们的服侍，而它们自己也被解脱束缚时，就更不会遭受不义了。因为，倘若有言语的恩赐，它们不会控诉天主，说祂不公义地使它们因未分享同样的复活而低人一等。因为对于本性不同的受造物，公义的存在者不会分配给它们同等的结局。而且，对于那些没有正义概念的受造物，也不可能抱怨不公义。也不能说，对于那将要复活的人而言，存在着不义，因为他由灵魂和身体组成，他无论对灵魂还是身体都没有遭受不义。任何心智健全的人都不会断言他的灵魂遭受了不义，因为这样说时，他不经意间也影射了当前的生命；因为，如果现在，当灵魂居住在会腐朽且会受苦的身体中时，它并未遭受不义，那么，当它与一个免于腐朽和受苦的身体结合而生活时，就更不可能遭受不义了。再者，身体也不遭受不义；因为，如果现在，当一个可腐朽之物与一个不可腐朽之物结合时，它并未遭受不义，那么显而易见，当一个不可腐朽之物与一个不可腐朽之物结合时，它也不会遭受不义。不，也没有人可以说，使一个已经消散的身体复活并重新聚合起来，是不相称于天主的工作：因为，如果较坏的事——即造一个会腐朽且会受苦的身体——对祂而言并不是不相称的，那么，较好得多的事——即造一个不会腐朽或受苦的身体——就更不是不相称的了。\par

% structure: p0022 source=h2 target=section reason=relative-heading-rank; base=h2

\section{第十一章 重述要点}

如果，那么，通过那种本质上在先的和从之而来的东西，所探讨的每一点都已得到证明，那么十分明显的是，已分解的身体的\OriginalTerm{复活}{resurrection}是\OriginalTerm{造物主}{the Creator}能够施行、愿意施行、且与祂相称的一项工程：因为通过这些思考，相反意见的虚妄已经显明，不信者所采取的立场之荒谬也已昭示。我为何还要提及它们彼此之间的对应，以及它们之间的关联呢？如果应当使用关联一词，仿佛它们因某种本性上的差异而分离；而不如说，天主所能的，祂也愿意；天主所愿意的，祂完全能够实现；并且这合乎那位愿意者的尊荣。论说真理是一件事，为真理辩护是另一回事，已在前述评注中充分说明，同样说明了它们在哪些方面彼此不同，以及何时、对付何人时各自有用；但或许没有理由，为了普遍的确定性，并由于已说的与余下部分的关联，我们不从这些相同点和与之相关之点重新开始。一种论证方式自然位居首位，另一种则紧随其后，清理道路，消除阻碍或敌对。论说真理，因为对所有人为确定和安全所必需，故在本性、次序或效用上皆居首位：在本性上，它提供关于主题的知识；在次序上，它在于它所告知的事物中并与它们同在；在效用上，它是那些了解它的人之确定与安全的保证。为真理辩护的论证，其本性和效力次之，因为驳斥谬误不如确立真理重要；在次序上次之，因为它施力于持有错误意见者，而错误意见是别样种子和退化的后产物。但是，尽管如此，它常被置于首位，有时更显有用，因为它预先消除并扫清某些人心中的不信，以及那些初来者的疑惑或错误意见。然而，两者都归于同一目的，因为驳斥谬误和确立真理都以虔诚为目标：诚然，它们并非绝对同一，而是前者，如我所言，对所有信者及关注真理和自身\OriginalTerm{救恩}{salvation}的人皆是必需；后者在某些场合、对某些人和处理某些事上证明更为有用。至此，重述要点以忆及已说过的话。现在我们必须转入我们所提议的，并从原因本身（按照它、并因它之故，第一人及其后裔得以受造，尽管他们并非以同样方式进入存在）、从所有人的共同本性（作为人）、以及进一步从造物主按各人生活的时间、根据各人规范其行为的准则对他们进行的审判——这审判无人怀疑会是\OriginalTerm{公义}{just}的——来证明关于复活的道理之真实性。\par

% structure: p0024 source=h2 target=section reason=relative-heading-rank; base=h2

\section{第十二章 由创造人的目的论证复活}

从原因出发的论证将显现出来，如果我们考虑人是否是被随意且徒然受造，或是为了某种目的；如果是为了某种目的，那么是否仅仅为了使他能生活并继续存在于他受造时的自然状态中，或是为了供另一位使用；如果是为了使用，那么是为了造物主自身的使用，还是为了某些属于祂、并被祂认为配得更大照顾的存在者。现在，如果我们以最一般的方式考虑这一点，我们发现一个有健全心智，并受理性判断驱动而行动的人，不会徒然做任何他有意图去做的事，而是或者为了自己的用处，或者为了他所关心的某个他人的用处，或者为了作品本身的缘故，出于对其生产的某种自然倾向和喜爱。例如（为明确含义，我们使用一个说明），一个人为自己使用建造房屋，但为牛、骆驼和其他他所需要的动物，他建造适合它们各自的棚舍；如果不是为了自己的使用，如果我们只看表面，然而如果我们看他所意图的目的，那也是为了自己的使用，但就直接对象而言，是出于对他所关心之物的关切。他也生育子女，不是为了自己的使用，也不是为了属于他的任何其他事物，而是为了使由他而出的人能够存在并尽可能长久地延续，这样通过子女和孙辈的继承，他对自己生命的终结得到安慰，并希望以这种方式将可朽的变为不朽的。这就是人的行为方式。但是天主既不会徒然造人，因为祂是智慧的，智慧的作品没有一样是徒然的；也不会为了自己的使用而造人，因为祂一无所缺。但对于一个绝对无所需的存在者，祂的工程中没有一个能为祂自己的使用增添什么。再者，祂也不为了祂所造的其他工程中的任何一个而造人。因为没有任何被赋予理性和判断的事物，已经受造或正在受造，是为了供另一个——不论比它更大或更小——所使用的，而是为了那受造物自身的生命和延续。因为理性不能发现任何用途，可被视为创造人的原因，因为不朽者无所需，为了存在无需人的帮助；而无理性的存在者天生处于服从状态，并为它们各自的目的而为人提供服务，但它们反过来并不被预定去使用人：因为这既非正当，也不曾是正当的，去降低那统治和主导者供低等者使用，或使理性服从非理性——后者不适于统治。因此，如果人既非无因、徒然受造（因为天主的工程没有一样是徒然的，至少就造物主的意图而言），也不是为了造物主自己或由祂而出的任何工程所用，那么十分清楚，尽管按照最初且更一般的观点，天主为了自己造人，并遵循在创造中显耀的善与智慧，然而按照更切近受造物的观点，祂为了受造者的生命而造人，这生命并非燃起片刻而后熄灭。因为我想，对于爬虫、飞鸟、鱼类，或更一般地说，一切无理性的受造物，天主已分配给它们这样的生命；但对于那些背负造物主自身形象、被赋予理解、并享有理性判断的，造物主已安排了永久的延续，好使他们认识自己的创造者，及其能力与技巧，并服从法律和公义，从而在拥有那些品质中度过其整个存在，免受苦难，他们曾在可腐朽的和属地的身体中勇敢地忍受了先前的生命。因为凡是为了他物而受造之物，当那其为之受造的目的不复存在时，它本身也理当不再存在，并且不会徒然继续存在，因为在天主的工程中，无用之物无容身之地；但那为了存在和活出与其本性自然相合的生命而受造之物，既然原因本身与其本性捆绑在一起，且仅在关联于存在本身时才被认出，就绝不可能容许任何原因彻底消灭其存在。但由于这一原因被看作在于永恒的存在中，那如此受造之物就必须被永远保存，行与其本性相合之事并经历之，它由两部分组成，每一部分贡献其所当有的，好使灵魂按其所造的本性中无改变地存在并保持，并履行其适当的职能（诸如主宰身体的冲动，并以适当的标准和尺度判断和衡量不时发生的事），而身体则按其本性向着其适当对象运动，经历分配给它的变化，而在其中（涉及年龄、外貌或尺寸），复活也是其一。因为复活是一种变化，是所有变化中的最后一个，并且是一种对那时仍存留者变得更好的变化。\par

% structure: p0026 source=h2 target=section reason=relative-heading-rank; base=h2

\section{第十三章 论证的延续}

我们对这些事深具信心，一如对那些已经成就的事；并且，反思我们自己的本性，我们满足于与贫乏和朽坏相连的生命，因为这适合我们现今的生存状态；同时，我们坚定不移地盼望在不朽中持续存在；我们这样做，并非毫无根据地建立在人的虚构之上，以虚假的盼望自欺，我们的信仰乃是基于一个最绝对可靠的保证——那位创造我们者的目的，依此目的，祂造了人类，具有不朽的灵魂与身体，并赋予他理智与内在的律法，为保守并维护祂所赐予的一切——这些恩赐与一个有理智的存在和理性的生活相称：因为我们深知，倘若祂无意让如此受造的物存续至永恒，祂就不会形成这样的存在，也不会赋予他属于永恒的一切。因此，倘若这位宇宙的造物主造人，是为了让他能分享理性的生活，并且让他成为祂的伟大以及显明在万物中的智慧的观赏者，使他能永远不断地深思这些；那么，依照他的创造者的目的，以及他所领受的本性，他被造的缘由就是他永远存续的保证，而这存续又是复活的保证，没有复活，人便无法持续存在。这样，根据所说的，显而易见，复活已经由人受造的缘由以及那造他者的目的清楚证明了。既然人进入这个世界的原因的本性如此，接下来要思考的就是自然紧随其后的、或按预定顺序的事情；而且，在我们探讨中，紧随受造的缘由之后的是如此受造的人的本性，以及造物主因公正审判而施于他们的那些人的本性，而所有这一切之后则是他们存在的终向。因此，在业已考察了排列在先的问题之后，我们现在必须进而思考人的本性。\par

% structure: p0028 source=h2 target=section reason=relative-heading-rank; base=h2

\section{第十四章 复活并非仅基于未来审判的事实}

由真理所组成的各项教义，或任何待检验的事，若要产生对所言之事的不可动摇的信心，其证明必须不是始于任何外在事物，也不是始于某些人现今或曾有的看法，而是始于关于此事之共通而自然的观念，或始于次要真理与首要真理的联系。因为问题要么涉及首要信念，那么只需回忆，以激发自然观念；要么涉及由首要事物自然引出者及其自然顺序。对于这些事，我们必须遵循次序，指出哪些是从首要真理或从先置者严格推演而来，如此我们既不会忘记真理或对它的确定性，也不会混淆那些由自然安置且彼此有别的事物，或打乱自然次序。因此，我认为，凡欲公平处理此主题，并愿意明智判断复活是否存在的人，首先应仔细思考有助于证明此事的论证力度及各论证所处位置——何者为第一、何者第二、何者第三、何者最后。在安排这些论证时，应将人受造的缘由置于首位——即创造者造人的目的；然后，按其适合程度，将此与如此受造者的本性联系起来；并非因为本性在次序上居次，而是因为我们无法同时判断两者，尽管它们相互具有最紧密的自然联系，且对当前主题具有同等效力。然而，尽管出自这些首要证据，且由创造工程引申出来的论证已清楚证实了复活，我们仍可从取自天主眷顾的论证中获得关于复活的信念——我意指按公义审判给予每人的赏罚，以及人类存在的终向。因为许多人在讨论复活这一事项时，将全部理由都仅仅基于第三个论证上，认为复活的缘由就是审判。但此一谬误极为明显，因为虽然所有死去的人都要复活，然而并非所有复活的人都要受审判：因为倘若只有公义审判才是复活的缘由，那么当然可将那些既未行恶也未行善者——即年幼孩童——排除在复活之外；但既然众人都要复活，包括那些夭折的婴孩和其他人，他们同样证明我们的结论：复活的发生，并非以审判为首要理由，而是由于天主形成人的目的，以及如此受造者的本性。\par

% structure: p0030 source=h2 target=section reason=relative-heading-rank; base=h2

\section{第十五章 由人的本性对复活的论证}

虽然从人的受造中可发现的原因本身足以证明，复活是随身体消亡而自然产生的后续结果，但或许仍应毫不退缩地引证前述论证中的每一种，并依前文所言，向那些自己无力辨识的人指出由最初真理衍生出的每一真理中的论证；首先且最重要的，是受造之人的本性，它将我们引向同一观念，并具有与复活证据同样的效力。因为，若全人类的本性普遍地由不朽的灵魂与在受造时与之相配的身体构成，且若天主既非单独将灵魂的本性，亦非单独将身体的本性，赋予这样的受造、这样的生命及整个存在过程，而是赋予二者结合而成的人，好使他们在度过今世的存在之后，得以由他们出生及在世时所由构成的相同元素，抵达一个共同终向；既然一个活物由二者形成，经历灵魂所经历的一切与身体所经历的一切，并执行和完成任何需要感觉或理性判断的事，则必然推断，这一切系列都必须指向某一个终向，好使它们全体，并藉由全体——即人的受造、人的本性、人的生命、人的作为与苦难、他的存在过程以及适合其本性的终向——都能和谐一致，并拥有相同共通的体验。然而，若整个存在者——无论是源于灵魂的事物，或是藉由身体所成就的事物——都拥有某种一体的和谐与共通的体验，那么这一切的终向也必是同一的。严格说来，若那终点所归属的存在者在构成上保持同一，终向便是同一的；而若构成该存在者之各部分的那些事物保持同一，该存在者便将全然相同。若已分解的各部分为着构成该存在者而重新结合，它们在各自特有的结合上便是相同的。而同一世人的构成，必然证明已死且已分解的身体将会复活；因为若无复活，同一的各部分既不能按本性彼此结合，同一世人的本性也无法重新构成。再者，若理智与理性被赋予人，是为了辨识那由理智领悟的事物，不仅辨识其存在，也辨识其赐予者的美善、智慧与公义，那么必然可得：既然理性判断所为之持续的那些事物得以持续，为这些事物赐予的判断也当持续。但除非那接受了它并维系着它的本性持续存在，否则此事不可能持续。然而，那接受了理智与理性的，是人，而非单独的灵魂。故此，由两部分构成的人，必须永远持续。但若他不复活，便不可能持续。因为倘若没有复活，人之作为人的本性便无法持续。而倘若人的本性不持续，灵魂之适配于身体的需要与其经验便归于徒然；身体之被束缚而不得获其渴求的，受制于灵魂的缰绳并如被嚼环引导，亦归于徒然；理智归于徒然，智慧归于徒然，遵守公义乃至实践一切美德归于徒然，法律的制定与执行亦归于徒然——一言以蔽之，人内的一切高贵事物或为人而存在的事物，甚或人的受造与本性本身，皆归于徒然。然而，若虚空全然被排除于天主的一切工作及祂所赐的一切恩惠之外，则不可避免的结论便是：在灵魂无尽的延续之余，身体亦将依其固有本性而永久存续。\par

% structure: p0032 source=h2 target=section reason=relative-heading-rank; base=h2

\section{第十六章 死亡与睡眠的类比，以及由此得出的复活论证}

谁也不要认为，我们把被死亡和腐朽中断的存在的存续称为生命是一件奇怪的事；相反，他应当考虑到，这个词并非只有一种含义，存续的尺度也不止一种，因为那些持续存在的事物的本性也不是单一的。因为，如果每一种持续存在的事物都按照其特有的本性而存续，那么对于那些完全不可腐朽且不死不灭的存在者，我们也找不到与我们的存续相同的存续，因为高等存在者的本性并不降低到低等存在者的水平；在人身上，寻求一种不变不迁的存续也是不恰当的；因为前者从起初受造便是不死不灭的，并仅凭其造物主的单纯意愿而无止境地继续存在；而人，就其灵魂而言，从起初便有一种不变的存续，但就身体而言，则通过变化而获得不死不灭。这就是复活教义的含义；我们着眼于此，既等待身体的解体，视之为一种充满匮乏与腐朽的生活的结局，又在此之后希望获得一种带有不死不灭的存续，既不将我们的死亡等同于无理性的动物的死亡，也不将人的存续等同于不死不灭者的存续，免得我们无意间以这种方式将人的本性与生命放在不适宜与之比较的事物之列。因此，如果在人的寿命方面似乎存在某种不均，也不应引起不满；我们也不应因为灵魂与身体肢体的分离以及身体各部分的解体中断了生命的连续性，就因此对复活感到绝望。因为，虽然感官和身体能力的松弛——这自然发生在睡眠中——当人们每隔相等的时间间隔入睡，并似乎又活过来时，似乎中断了感觉生命，但我们并不拒绝称之为生命；而且，我想，正因如此，有人称睡眠为死亡的兄弟，不是因为他们出自同样的祖先和父亲，而是因为死者和睡眠者处于相似的状态，至少就静止不动、完全感觉不到现在或过去，甚至感觉不到存在本身和自己的生命而言。因此，如果对于从出生直到解体，充满此种不均，并被我们前面所提及的一切事物所中断的人的生命，我们都不拒绝称之为生命，那么对于继解体之后的、包含复活在内的生命，我们也不应感到绝望，尽管它在一段时间内因灵魂与身体的分离而中断。\par

% structure: p0034 source=h2 target=section reason=relative-heading-rank; base=h2

\section{第十七章 我们现在在人身上能追溯的一系列变化使复活成为可能}

因为人的这种本性，从一开始便被定下了不均，并且按照其造物主的旨意，拥有一种不均的生命和存续，时而为睡眠所中断，时而为死亡所中断，且为每个生命阶段所发生的变化所中断，而那些紧随最初阶段之后的变化事先并未清楚显现。若非经验教导，谁会相信，在那各部分都同样柔软的种子中，竟蕴藏着如此多样且数量庞大的巨大能力，或者团块——即骨骼、神经、软骨，还有肌肉、肉、肠子以及身体的其他部分——它们以这种方式产生并变得坚实？因为，在仍然湿润的种子中，这类东西无一可见；甚至在婴儿身上，那些属于成年人的特征也无一显现；在成年时期，也不显现那些属于已过壮年的人的特征；在这些人身上，也不显现那些属于已衰老之人的特征。但是，尽管我所列举的一些事例完全没有显示、另一些也只模糊地显示了人类本性中的自然顺序和随之而来的变化，然而凡是未被恶习或怠惰蒙蔽了对这些事理的判断的人，都知道：必须先有种子被放置；当这种子按每一肢体和部分完全组织成形，子嗣得见光明时，便有属于生命最初阶段的生长，以及伴随生长的成熟；成熟之后，是身体能力的衰退，直至老年；然后，当身体衰败时，便是它的解体。因此，在这件事上，尽管种子上面并未铭刻人的生命或形式，生命本身也未铭刻分解为原初元素的命运；自然事件的承续却使得那些从现象本身看来并不可信的事物变得可信；那么，理性就更其如此——它从自然序列中追寻真理，为相信复活提供了根据，因为在确立真理上，理性比经验更稳妥、更可靠。\par

% structure: p0036 source=h2 target=section reason=relative-heading-rank; base=h2

\section{第十八章 审判必须同时涉及灵魂与身体：因此将有复活}

我刚提出的、为证明\OriginalTerm{复活}{resurrection}真理而供检验的那些论证，都属于同一类，因为它们都从同一点出发；因为它们的出发点是创造中首批人类的起源。但是，其中一些论证的力量来自它们所源的出发点本身，另一些则根据人的本性和生命，从天主对我们的眷顾中获得可信性；因为人之所以形成的依据和缘由，与人的本性紧密相连，其力量源于创造；而源自\OriginalTerm{正义}{rectitude}的论证——它表明天主按照人生活的好坏来审判人——的力量则来自于他们存在的终结：人的存在基于前者的理由，但其状态更依赖天主的眷顾。既然我已尽我所能证明了在先的事情，那么也宜用随后的事情来证明我们的论点——我指的是按照公义的审判每个人应得的赏罚，以及人类存在的终极原因；其中，我将本性上领先的放在首位，首先探究有关审判的论证：只预先提出一件事，出于对我们面前之事所关乎的原则和秩序的关切——即，凡承认天主是这宇宙的创造者的人，若想坚持自己的原则，就必须将一切受造物的保存和眷顾归于祂的智慧和正义；并持这样的观点：认为天地间没有一件事物不受守护和\OriginalTerm{眷顾}{providence}，相反，创造者的关注遍及一切，无论是无形有形、是微是大；因为所有受造物都需要创造者的关注，每一物更是按照其本性和受造的目的而需要：虽然我认为现在逐一列举、区分各种情况、或详细说明什么适合于每一本性，将是白费力气。无论如何，人就我们现在所要谈论的而言，因处于匮乏之中，就需要食物；因是可朽的，就需要后代；因是有理性的，就需要审判的过程。但如果这些事物都按本性属于人，他需要食物以维持生命，需要后代以延续种族，需要审判以使食物和后代都能合乎法则，那么自然可推知，既然食物和后代都涉及两者结合的整体，审判也必须指向这整体（我所说的两者结合的整体是指由灵魂和身体组成的人），并且这样的人要对他的所有行为负责，并因此而获得赏或罚。现在，如果公义的审判要将赏罚归于行为所由行的整体；并且，若要让灵魂单独承受与身体结合时所行行为的报应，是不恰当的（因为灵魂本身对涉及快乐或身体的食物与修养的过愆并没有倾向），而让身体单独承受也同样不当（因为身体本身不能分辨法律与正义），但由这两者组成的人却要为他所行的每件事受审；而理性既未发现这在此生中发生（因为按功过行赏在今生中无处可寻，因为许多无神论者和行尽一切邪恶的人终其一生未遇灾祸，相反，那些在一切德行上明显生活典范的人却生活在痛苦、侮辱、诽谤和残害以及种种苦难之中），也未发现在死后发生（因为整体已不再存在，灵魂与身体分离，身体本身又分解成其组成的元素，不再保留其先前的结构或形态，更遑论对其行为的记忆）：这一切的结果对每一个人都是非常明显的——即，用宗徒的话说，\BibleQuote{这败坏（且可分解的）必须穿上不朽坏的，}\BibleRef{格前 15:54}好使那些已死的人，藉着复活重获生命，被分离且完全分解的部分得以重新结合，每个人都能按公义，领受他藉身体所行的，无论善恶。\par

% structure: p0038 source=h2 target=section reason=relative-heading-rank; base=h2

\section{第19章 若无复活，人之处境较禽兽更为不利}

答复那些承认天主的上智安排、并且接受与我们相同原则的人——但他们却不知何故偏离了自己的承认，我们可以使用诸如先前所提出的论证，以及更多于此的论证，倘若有人愿意扩展那仅被简要而粗略提及的论述。然而，对于在基本真理上与我们意见相左的人，或许最好提出另一个先于这些的原则，与他们一同怀疑他们意见所涉之事，并以此方式与他们一同考察——人类的生命及其整个存在过程是否被忽略，似乎有一层浓厚的黑暗倾注于世，将人及其行为隐藏在无知与沉默之中；抑或，更有把握的看法是，造物主亲自掌管祂自己所造的万物，监察一切存在或生成的事物，是行为与意念的审判者。因为，倘若对人的行为完全不加审判，那么人比起无理性的受造物便毫无优越之处，反而处境更劣，因为他们在克制情欲、关心虔敬、公义以及其他德性；而禽兽般的生活方式便成为最佳，德性成为荒谬，审判的威胁成为笑柄，放纵各种享乐成为至善，所有这些人的共同决定和唯一准则，便是那为放纵和淫乱之人所钟爱的格言：\BibleQuote{我们吃喝吧！因为明天就要死了。}\BibleRef{格前 15:32} 因为这种生命的结局并非如某些人所想象的享乐，而是完全的麻木。然而，倘若人类的造物主对自己的工程有所关心，并且善生与恶生的人之间确有区别，那么这种区别，要么是在今生，当人行善或作恶时仍在世的时候；要么是在死后，当人处于分离与消散的状态。但是，按照这两种设想，我们都无法找到公正的审判发生；因为今生，善人并未获得德性的赏报，恶人亦未得恶行的报应。我暂且不论这一点，即只要我们现在所拥有的本性得以保存，道德本性便无法承受与众多或更严重过错相称的惩罚。因为盗贼、统治者或暴君，曾不义地杀害无数众人，不可能以一次死亡抵偿这些罪行；至于那对天主毫无正确认识、却生活在一切暴行与亵渎中的人，他藐视神圣之事，违犯法律，同样侵害男孩与妇女，不义地摧毁城市，焚烧房屋连同居民，蹂躏国土，同时灭绝城邑居民与百姓，甚至整个民族——他如何能在这可朽坏的身体内，承受与这些罪行相称的惩罚呢？因为死亡阻止了应得的惩罚，而且这有死的本性不足以抵偿他的任何一件恶行。因此，得证无论在今生还是在死后，都没有按人所应得的审判。\par

% structure: p0040 source=h2 target=section reason=relative-heading-rank; base=h2

\section{第二十章 人必须在来世具有身体和灵魂，审判才得公正}

因为死亡要么是生命的彻底灭绝，灵魂与身体一同消散败坏；要么是灵魂自身存留，不散不灭、不腐败，而身体则腐蚀消散，不再保留任何对过去行为的记忆，也不再感知与灵魂相关的经历。倘若人类生命要彻底灭绝，那么显然，不会有人关心那些不再活着的人，不会有针对那些善生或恶生之人的审判；反而，一切放荡不羁的生活，以及由此而来的众多荒谬，乃至这一不法行为的顶点——\OriginalTerm{无神论}{atheism}，都将再次侵袭我们。然而，倘若身体会腐败，而分解的每一部分都归于同类的元素，但灵魂却自身存留，不死不灭，那么按照这种设想，灵魂亦不会受到任何审判，因为这将缺乏公义：因为从天主而来、且源于天主的审判若缺乏公义，是不可想象的事。然而，如果那曾行义或不义的人本身没有保存下来，审判\Emph{确实}缺少公义：因为那在生活中行了审判所依据诸事的，是人，不是灵魂本身。一言以蔽之，这种观点无论如何不符合公义。\par

% structure: p0042 source=h2 target=section reason=relative-heading-rank; base=h2

\section{第二十一章 论证的继续}

因为如果善行受到赏报，身体显然会被亏待，因为它曾与灵魂一同分担了行善的辛劳，却未分享善行的赏报；而且，虽然灵魂常因身体的需求和缺乏而被宽恕某些过失，但身体本身却被剥夺了分享所行善行的一切份子，这善行乃它生前协助承担的辛劳。同样，如果过失受到审判，灵魂也得不到公正的对待，倘若它独自偿还因受身体引诱而犯下的过失的惩罚——灵魂被身体拖向身体的欲望和动向，有时被攫住强行带走，有时以某种极其暴力的方式被吸引，有时则出于善意及对身体保存的关照而迁就于它。怎么可能不是不公的呢：灵魂就其本性对放纵、暴力、贪婪、不义以及由此产生的不义行为这类事，本无任何欲望、任何动向、任何冲动，却独自为此受审判？因为如果这类恶事大多来自人未能控制那些引诱他们的\OriginalTerm{偏情}{passions}，而他们又是受身体的需求与缺乏以及身体所需的照顾与留意所引诱（因为这些正是获取一切财产、尤其是使用财产的动机，而且也是婚姻及人生一切行动的动机，在这些事上并与此相关，过错与无过得以显明），那么那些事，身体是首先感知的，并为了身体所要的东西将灵魂拖向同情和参与行动，而欲望与快乐，以及恐惧与忧愁，其中任何超出适当界限的都要受审，这些竟由身体发动，然而由此产生的罪以及对这些罪的惩罚竟只落在灵魂身上，而灵魂既不需要这类事物，自身也根本不渴望、不畏惧、不忍受任何人惯常遭受之事——这又怎么可能公正呢？但即使我们认为这些情感不只属于身体，而是属于人——这样说应是正确的，因为人的生命是一个整体，尽管由两者组成——然而若我们仅观其固有本性，我们绝不会断言这些事属于灵魂。因为如果灵魂完全无需食物，它就绝不会渴望那些对它的生存毫无所需的事物；它也不会对那些完全不适合它使用的事物产生任何冲动；同样，它也不会因缺钱或缺少其他财产而悲伤，因为这些东西并不适合它。而且，如果它超越败坏，它就无所畏惧，不惧任何能毁灭它自身的事物：它不惧怕饥荒、疾病、残损、瑕疵、烈火、刀剑，因为它不会从其中任何一样遭受伤害或痛苦，因为身体或身体的能力根本不能触及它。但若将偏情归于灵魂，视之为特别属于它的，这是荒谬的，那么仅仅将源于偏情的罪及其应受之惩罚加于灵魂身上，便是极度不公，且不配天主的审判。\par

% structure: p0044 source=h2 target=section reason=relative-heading-rank; base=h2

\section{继续论证}

除了已说过的，这岂不荒谬：我们甚至不能设想\OriginalTerm{德行}{virtue}与\OriginalTerm{恶习}{vice}单独存在于灵魂中（因为我们承认德行是人的德行，同样，恶习作为其对立面，也不属于与身体分离、独立存在的灵魂），然而对这些的赏或罚却单单归于灵魂？当灵魂没有对死亡、创伤、残损、损失、凌辱，或与之相连的疼痛，或由此而来的苦难的恐惧时，有谁能设想\OriginalTerm{勇敢或刚毅}{courage or fortitude}单独存在于灵魂中呢？当没有欲望将它拉向食物、性交，或其他快乐与享受，也没有任何其他事物从内部引诱它或从外部激发它时，我们又该怎么说\OriginalTerm{自制与节制}{self-control and temperance}呢？当事物中没有任何可行或不可行的提议，也没有要选择或避免的事物，或者不如说，当它内根本没有任何向行动的行动或自然冲动时，\OriginalTerm{实践智慧}{practical wisdom}又是什么呢？\OriginalTerm{公平}{equity}又如何能以任何意义成为灵魂的属性，无论涉及彼此还是任何别的事物，无论是同类的还是不同类的，当它们不能从任何来源、以任何方法、借任何途径，按照功绩或类比给予那公平的事物（除了对天主的尊敬外），而且，对使用自己的东西或禁戒别人的东西，也没有任何冲动或动向，因为对依本性之事的使用或禁戒，是相对于那些被构成能使用它们的人而言的，而灵魂既无所需求，也未被构造成能使用任何事物或任何单一事物，因此所谓各部分独立行动，不能发现于如此构成的灵魂中？\par

% structure: p0046 source=h2 target=section reason=relative-heading-rank; base=h2

\section{继续论证}

但最不合理之事莫过于此：为人类制定适当的法律，然后却单单将守法或违法行为的报偿归于他们的灵魂。因为，如果领受法律的人也理当承受违法的报偿，而且领受法律的是人，而非灵魂本身，那么人也必须承受所犯罪过的报偿，而非灵魂本身，因为天主并未命令灵魂戒避那些与它们无关的事，例如奸淫、杀人、偷窃、抢劫、不孝敬父母，以及一切倾向于伤害和损害邻人的贪欲。因为那\BibleQuote{应孝敬你的父亲和母亲}的命令并非单单针对灵魂，因为这些称谓对它们并不适用，因为灵魂并不产生灵魂，从而配得父亲或母亲的称呼，而是人产生人；同样，那\BibleQuote{不可奸淫}的命令，也绝不能恰当地对灵魂发出，甚至在这样的语境中根本不可设想，因为在它们之中并没有男女之别，也没有任何性交的能力或欲望；而没有欲望，便不会有交合；而根本就没有交合，便不会有合法的交合，即婚姻；而没有合法的交合，也就不可能有对他人妻子的非法欲望或交合，即奸淫。再者，禁止偷窃，或禁止贪多务得的命令，对灵魂也不适用，因为它们并不需要那些东西——人出于本性的匮乏或缺乏，便习惯偷窃或抢劫的东西，例如金、银、牲畜，或其它适宜作食物、遮蔽物或用具的物品；因为对不死不灭的本性而言，贫穷者视为有用的一切，都是无用的。但更详尽的讨论，就留给那些希望更精确地考察每一点，或更激烈地与对手争辩的人吧。然而，既然方才所说的，以及与此协同保证复活的论述，对我们已经足够，再继续详述它们就不再合宜了；因为我们的目标并非毫无遗漏地陈明一切可说的，而是以概括的方式向聚集的听众指出对于复活所当持守的看法，并使有关这一问题的论证适应在场之人的接受能力。\par

% structure: p0048 source=h2 target=section reason=relative-heading-rank; base=h2

\section{第二十四章 从人的最终目的论证复活}

所提出来考虑的各点既已在一定程度上得到考察，余下的就是考察从目的或\OriginalTerm{目的因}{final cause}而来的论证，这一论证其实在已有的论述中已经浮现，只需要稍加注意和进一步的讨论，便能避免留下我们未提及任何简要提及之事的印象，从而间接损害主题或起初所作的主题划分。为此，为了在场的各位，以及可能关注这一主题的其他人，或许正好可以指明：凡是本性所构成的事物，以及艺术所造就的事物，都必然有其自身所特有的目的，正如全人类的常识所教导我们的，并为我们眼前的事物所证实的。我们岂不看见，农人有农人的目的，医生有医生的目的；再者，从土里生长出来的事物有它们的目的，而靠土得养、依循一定的自然秩序产生的动物，也有它们的目的吗？如果这是明显的，并且自然与人为的能力，以及由这些能力所产生的行动，都必须伴随着符合本性的目的，那么绝对必要的是，人的目的，由于它属于一种特殊的本性，应当与其它事物的目的区别开来；因为，我们不可认为缺乏理性判断的存在者，与那些行事受内在法律和理性规范、过理智生活并持守正义的存在者，具有相同的目的。因此，免于痛苦不能是后者的恰当目的，因为这目的他们与全无感觉的存在者共有；这目的也不可能在于享受那些滋养或愉悦身体的事物，或在于丰盈的快乐；否则，一种如同野兽般的生活必占据首位，而那受德行规范的生活却失去了目的因。因为，我认为，这样的目的属于野兽和牲畜，而不属于拥有不朽灵魂和理性判断的人。\par

% structure: p0050 source=h2 target=section reason=relative-heading-rank; base=h2

\section{第二十五章 论证续完}

也不是灵魂与身体分离后的幸福：因为我们探讨的不是构成人的两个部分中任何一个的生命或\OriginalTerm{目的因}{final cause}，而是由二者结合而成的存在者；因为每个分有此现世存在的人都是如此，并且必然有为此生命所定的某种适当目的。但如果这是两个部分结合的目的，而这目的既不能在他们仍生活于此现有存在状态时，因前面提到的众多原因而得以发现，也不能在灵魂处于分离状态时发现，因为当身体分解，甚至全然消散时，即使灵魂独自存续，也不能说这个人仍然存在——那么，绝对必然的是，人之存在的目的应显现在二者的某种重新结合中，并且是同一个生命存在的重新结合。既然这出于必然，那么必定有死去的、甚至完全分解的身体的复活，同样的人必须被重新形成，因为自然的法则并非绝对地为任何人设定目的，也不是为随便什么人，而是为那些经历过先前生活的同样的人；但除非同样的身体归还给同样的灵魂，否则同样的人是不可能重新构成的。而同样的灵魂获得同样的身体，在任何其他方式下都是不可能的，只有通过复活才可能；因为如果这事发生，那么适合人之本性的目的也随之而来。我们说，理智生命和理性判断的\OriginalTerm{目的因}{final cause}，就是不断专注于那些自然理性首要且主要适应的对象，并不断在对\Emph{自有者}及其谕令的瞻想中喜乐，尽管大多数人因为过分热情和强烈地受到下面事物的影响，度过一生而未达到这一目的。然而，众多未能达成其当有之目的的人，并不会使共同的命运落空，因为审查是针对个人的，而对善度或虚度一生的奖赏或惩罚，是与各人的功过成比例的。\par

\SourceRecord{Translated by B.P. Pratten.；Ante-Nicene Fathers；Vol. 2.；Edited by Alexander Roberts, James Donaldson, and A. Cleveland Coxe.；Buffalo, NY: Christian Literature Publishing Co.,；1885.；<http://www.newadvent.org/fathers/0206.htm>.}
